\section{材料中的辽：制度名称之外的人与地点}

道宗在位的1056--1101年，是辽的制度、宗教和东北网络仍在运转的一段长时段；编年材料以连续在位纪年保存这段行政时间，而不是把它误写成无事发生的空白。要理解这类连续性，《辽史》的百官、营卫、食货等志只能给出后设整理的名目。契丹大字、小字题记、汉文墓志、寺院碑刻与城址，才可能把官号、家族、功德、道路和地方活动重新放回具体地点。

然而，材料的互补不等于材料会说同一种话。一件墓志说明的是追赠、亲属与精英自我陈述，不能据此统计某地的官员或族群；一座寺院的重修题记能显示施主和文字实践，也不等于国家意志或全体民众的信仰。尤其契丹文字的断片、释读与旧藏来源，必须先确认同件、同年与转写方案。辽的多中心秩序应由这些差异来理解，而不是用“北面”“南面”两个静态盒子代替。

东北并非只在1114年女真起兵时才进入历史。它是营卫、贸易、移民、渤海旧地城市与宗教网络长期交会的空间。辽朝对不同地区的行政、征发和军事控制不可能同密；从上京到南京的都城与交通节点，也不能被压缩为一条自北向南的权力链。正因为后期证据零散，书写道宗末年时尤其要分开国家法令、地方实践和墓葬所呈现的身份。

\section{宗教、文字与国家能力的边界}

辽代佛寺、造像、经幢和墓葬说明宗教资源能把家族、城市和官府联结起来，但不应从图像题材反推出统治者的单一意图。与其问辽是否被某一制度或某一“文化”完全定义，不如追问在1056--1101年的连续纪年中，谁以何种语言立碑、捐施、任官、迁徙或承担劳役。现存材料并不均衡，因此这种问题首先要求承认看不见的区域和人群。

这种限定也改变我们理解1114年以后崩解的方式。女真军事压力固然重要，辽朝原有的文书、寺院、家族与城市网络却不会因皇室覆亡而同时停止。将政权终结、行政接收、宗教延续和居民迁徙分别处理，才能使东北和燕云的社会史不沦为战争年表的附录。

\subsection{女真使者的请求：东北关系并非辽宋两家的余波}
\begin{event}{\eventanchor{SYD-991-LIAO-019}{女真请伐契丹，宋廷不许}}{淳化二年（九九一年）}{宋、辽与女真诸部}{宋廷回覆可核；请战者的范围、辽境内反应及女真内部政治不能由一条本纪概括}
九百九十一年女真使者向宋廷提出伐辽的请求时，东北并不是辽宋外交版图之外的空处。宋廷选择不许，使者所带来的地方竞争、边市信息和可能的军事联络没有立即化作一场宋辽战争；这个决定同样显示，女真行动者能够把自己的利益带到两大政权之间，而不只是被谁“控制”的对象。

《宋史·太宗纪》所留的是宋廷收到请求后的处置，《辽史》与东北考古、后出的区域材料所见又不相同。它们可以让我们确认一次交涉存在，不能替各部划出固定疆界或为其指定统一目标。后来东北权力关系的变化，正须从这些不对称的接触开始追问。
\end{event}